\section{First-order logic}
\label{sec:first-order-logic}
\slidesrc{FOL: 1--21; Proof theory (FOL): 1--4}

Propositional logic treats ``all monitors are ready'' as one indivisible block, true or
false. It cannot look \emph{inside} the sentence and see that it talks about \emph{all}
objects of a certain kind. \textbf{First-order logic} (FOL, also called \emph{predicate
logic}) can: it has names for objects, functions, relations between objects and, above
all, \textbf{quantifiers} (``for all'', ``there exists''). It is the logic behind Prolog
and behind the beliefs of BDI agents.

This chapter covers, in the order of the slides:
\begin{enumerate}[nosep]
  \item why propositional logic is not enough;
  \item the syntax of FOL: terms (constants, variables, functions) and predicates;
  \item quantifiers ($\forall$, $\exists$, $\exists_1$) and how to translate English into FOL;
  \item decidability: FOL has no truth-table method, and no method at all that always works;
  \item proof theory for FOL: Frege/Hilbert-style axioms, soundness, the sequent calculus LK.
\end{enumerate}

\subsection{Why not propositional logic?}
\label{sec:why-fol}
\slidesrc{FOL: 2--3}

Consider the argument:
\begin{center}
\begin{tabular}{l}
all monitors are ready \\
X12 is a monitor \\
\hline
therefore X12 is ready
\end{tabular}
\end{center}
Intuitively it is \emph{sound}: if we accept that the two premises (above the line) are
true, we must accept that the conclusion is true too. Each of the three statements is a
proposition (true or false), so we could try propositional logic:
\begin{center}
$p$ = ``all monitors are ready'', \quad $q$ = ``X12 is a monitor'', \quad $r$ = ``X12 is ready''.
\end{center}
But then the argument becomes $p,\ q \;/\therefore\; r$, which is \textbf{not} valid: the
valuation $p = q = \true$, $r = \false$ makes the premises true and the conclusion false.
The sense of the argument is lost, because the connection between ``monitors'',
``X12'' and ``ready'' is hidden inside the atomic propositions.

\begin{nbintuition}[Propositional logic is too coarse-grained]
Propositional logic sees sentences as black boxes. The validity of the monitor argument
depends on what is \emph{inside} the boxes: the same individual (X12), the same property
(being a monitor, being ready), and the word ``all''. We need a logic that can name
individuals, express properties and relations, and quantify over all individuals:
first-order logic. In FOL the argument becomes
\[
\forall x\,(\mathit{Monitor}(x) \limp \mathit{Ready}(x)),\quad \mathit{Monitor}(X12)
\;/\therefore\; \mathit{Ready}(X12)
\]
and it is valid (Section~\ref{sec:proof-theory-fol}).
\end{nbintuition}

\begin{nbwhy}
The slides come from a Software Engineering course, where FOL is used for \textbf{formal
specification} (e.g.\ the Z notation): ``every reactor has a temperature between 100 and
1000'' is a requirement about \emph{all} reactors. In AI the same language describes
knowledge bases (``every employee has a manager''), database queries (SQL is essentially
FOL), and the beliefs of agents.
\end{nbwhy}

\subsection{Syntax of FOL: terms and predicates}
\label{sec:fol-syntax}
\slidesrc{FOL: 4--13}

First-order logic is a generalisation of propositional logic with much greater expressive
power. Its language has two layers:
\begin{itemize}[nosep]
  \item \textbf{terms} talk about \emph{objects} (numbers, people, monitors);
  \item \textbf{predicates} (atoms) talk about \emph{facts} on those objects, and are true
        or false.
\end{itemize}
Formulas are then built from atoms with the connectives of propositional logic and with
quantifiers (Section~\ref{sec:quantifiers}).

\subsubsection{Terms: constants, variables, functions}
\slidesrc{FOL: 5--10}

\begin{nbdef}{Term (slide 5)}
The basic components of FOL are \textbf{terms}. A term denotes some \emph{object}, other
than true or false.
\end{nbdef}

The slides use a \emph{typed} version of FOL (as in the Z specification language): every
object belongs to a set, its \emph{type} ($\N$, $\Z$, Person, Monitor, \ldots).

\begin{nbdef}{Constant and variable (slide 6)}
\begin{itemize}[nosep]
  \item A \textbf{constant} of type $T$ is a name that denotes some \emph{particular}
        object in the set $T$. Examples: $8$ (the number eight, in $\N$ and $\Z$);
        $\mathit{mike}$ (the individual Michael Wooldridge); $S$ after ``let $S$ be the
        set $\{1,2,3\}$'' (by saying what it denotes we have given it an
        \textbf{interpretation}).
  \item A \textbf{variable} of type $T$ is a name that can denote \emph{any} value in the
        set $T$. A variable of type $\N$ can stand for any natural number.
\end{itemize}
\end{nbdef}

\begin{nbdef}{Functional term (slides 7--8)}
FOL has a set of \textbf{function symbols}; each denotes a function. Each function symbol
has an \textbf{arity} $n \in \N$ (its number of arguments), an argument type for each
argument and a return type. If $f$ is a function symbol of type $T$, with arity $n$ and
argument types $T_1, \ldots, T_n$, and $\tau_1, \ldots, \tau_n$ are terms of types
$T_1, \ldots, T_n$, then
\[ f(\tau_1, \ldots, \tau_n) \]
is a \textbf{functional term} (of type $T$).
\end{nbdef}

This is just like a function in a programming language: the slides' Modula-2 example
\code{PROCEDURE f(a1:T1; ... an:Tn) : T;} declares a function with $n$ typed arguments and
return type $T$.

\begin{nbexample}{the function \emph{plus} (slide 9)}
Let $\mathit{plus}$ take two natural numbers and return their sum. Then:
\begin{center}
\begin{tabular}{@{}ll@{}}
\toprule
\textbf{Acceptable terms} & \textbf{Not acceptable} \\
\midrule
$\mathit{plus}(2, 3)$ & $\mathit{plus}(-1, 0)$ \quad ($-1 \notin \N$: wrong type) \\
$\mathit{plus}(0, 1)$ & $\mathit{plus}(0.1, 2)$ \quad ($0.1 \notin \N$: wrong type) \\
$\mathit{plus}(\mathit{plus}(1, 2), 4)$ & $\mathit{plus}(1)$ \quad (wrong arity; added) \\
$\mathit{plus}(\mathit{plus}(\mathit{plus}(0, 1), 2), 4)$ & \\
\bottomrule
\end{tabular}
\end{center}
Terms can be nested: the argument of a function can itself be a functional term.
\end{nbexample}

Writing $2 + 3$ instead of $\mathit{plus}(2,3)$ is merely a convention (infix notation);
for logic, $+, -, /, *, \sqrt{\ }, \sin, \cos, \ldots$ are all function symbols. For
instance $(x + 3) * \sin 90$ is the term $\mathit{times}(\mathit{plus}(x, 3),
\sin(90))$ (slide 10).

\begin{center}
\begin{forest}
for tree={draw, rounded corners, font=\small, l sep=4mm, s sep=6mm, edge={thick}}
[$\mathit{times}$, fill=nbpurple!12
  [$\mathit{plus}$, fill=nbpurple!12 [$x$, fill=nborange!15] [$3$, fill=nbblue!10]]
  [$\sin$, fill=nbpurple!12 [$90$, fill=nbblue!10]]]
\end{forest}
\captionof{figure}{The term $(x+3)*\sin 90$ as a tree. Purple: function symbols; blue:
constants; orange: variables. Leaves are always constants or variables.}
\end{center}

\subsubsection{Predicates}
\slidesrc{FOL: 11--13}

\begin{nbdef}{Predicate (atom) (slide 11)}
FOL contains a stock of \textbf{predicate symbols}, standing for relationships between
objects. Each has an arity and a type for each argument. If $P$ is a predicate symbol of
arity $n$ with argument types $T_1, \ldots, T_n$, and $\tau_1, \ldots, \tau_n$ are terms of
those types, then
\[ P(\tau_1, \ldots, \tau_n) \]
is a \textbf{predicate} (also called an \textbf{atom} or atomic formula), which will be
either true or false under some interpretation.
\end{nbdef}

\begin{nbexample}{the predicate \emph{gt} (slide 12)}
Let $\mathit{gt}$ (``greater than'') take two natural numbers. Then $\mathit{gt}(4,3)$ is a
predicate that evaluates to \emph{true}, $\mathit{gt}(3,4)$ is a predicate that evaluates to
\emph{false}, while $\mathit{gt}(-1, 2)$ is not a predicate at all ($-1$ has the wrong
type). The usual mathematical relations $>, <, =, \geq, \leq, \neq, \in, \subset,
\subseteq$ are predicate symbols written infix. Predicates combine with the connectives:
$(2 > 3) \land (6 = 7) \lor (\sqrt 4 = 2)$.
\end{nbexample}

\begin{nbwarning}[Terms vs.\ predicates: the most common confusion]
A \textbf{term} names an \emph{object}; it is never true or false. A \textbf{predicate}
states a \emph{fact}; it is always true or false (under an interpretation).
$\mathit{plus}(2,3)$ is a term (it denotes 5); $\mathit{gt}(\mathit{plus}(2,3), 4)$ is a
predicate (it is true). Writing $\mathit{gt}(4,3) \land 5$ is nonsense: you cannot ``and''
a number. Likewise $\mathit{mother}(\mathit{ann})$ (a function: \emph{the mother of Ann}, a
person) differs from $\mathit{mother}(\mathit{mary}, \mathit{ann})$ (a predicate:
\emph{Mary is the mother of Ann}, a fact).
\end{nbwarning}

Some vocabulary (slide 13):
\begin{itemize}[nosep]
  \item a predicate with variables, e.g.\ $\mathit{gt}(x, 3)$, cannot be evaluated until we
        know what $x$ stands for (an interpretation of the variables);
  \item a predicate that contains \textbf{no variables} is a \textbf{proposition}
        (e.g.\ $\mathit{gt}(4,3)$);
  \item predicates of arity 1 are called \textbf{properties}, for example
        $\mathit{Man}(x)$, $\mathit{Mortal}(x)$,
        $\mathit{Malfunctioning}(x)$, all with one argument;
  \item predicates of arity 0 (no arguments) are called \textbf{primitive propositions}:
        they are exactly the atomic propositions of propositional logic. So propositional
        logic is the special case of FOL where all predicates have arity 0.
\end{itemize}

\begin{nbwhiteboard}[label=wb:terms-atoms]{Terms and atoms}{written after the FOL lectures}
\begin{center}
\renewcommand{\arraystretch}{1.25}
\begin{tabularx}{\linewidth}{@{}c X@{}}
\toprule
\textbf{Rule} & \textbf{Definition (as on the board)} \\
\midrule
1 & A \textbf{variable} is a term. \\
2 & A \textbf{constant} is a term. \\
3 & If $f$ is a \textbf{function symbol} with arity $m$, and $T_1, \ldots, T_m$ are terms,
    then $f(T_1, \ldots, T_m)$ is a term. \\
4 & If $p$ is a \textbf{predicate symbol} with arity $m$, and $T_1, \ldots, T_m$ are
    terms, then $p(T_1, \ldots, T_m)$ is an \textbf{atom}. \\
\bottomrule
\end{tabularx}
\end{center}
\begin{center}
\begin{tikzpicture}[font=\small, >=Stealth,
  b/.style={draw, rounded corners, thick, align=center, minimum height=0.8cm}]
  \node[b, draw=nborange, fill=nborange!10] (v) at (0,1.1) {variable\\[-2pt]\footnotesize $X$, $Y$};
  \node[b, draw=nbblue, fill=nbblue!10] (c) at (0,-0.4) {constant\\[-2pt]\footnotesize $\mathit{socrates}$, $8$};
  \node[b, draw=nbpurple, fill=nbpurple!10, minimum width=2.3cm] (t) at (3.6,0.35) {\textbf{TERM}\\[-2pt]\footnotesize denotes an object};
  \node[b, draw=nbgreen, fill=nbgreen!10, minimum width=2.3cm] (a) at (9.4,0.35) {\textbf{ATOM}\\[-2pt]\footnotesize true or false};
  \draw[->, thick] (v) -- node[above, sloped, font=\scriptsize]{rule 1} (t);
  \draw[->, thick] (c) -- node[below, sloped, font=\scriptsize]{rule 2} (t);
  \draw[->, thick, nbpurple] (t.north) .. controls +(0.4,1.3) and +(-0.4,1.3) .. node[above, font=\scriptsize, align=center]{rule 3: $f(T_1,\ldots,T_m)$} (t.north east);
  \draw[->, thick, nbgreen!70!black] (t) -- node[above=2pt, font=\scriptsize]{rule 4: $p(T_1,\ldots,T_m)$} (a);
\end{tikzpicture}
\end{center}
\wbphoto{sections/first-order-logic/fol-syntax/assets/whiteboard-6-1.jpg}
The definition is \textbf{recursive}: rule 3 builds terms out of terms, so terms can be
nested as deeply as we want ($f(g(X), c)$). Rule 4 is applied \emph{once}, on top: an atom
contains terms, but an atom is never an argument of a function or of another predicate.
Example with $X$ a variable, $\mathit{socrates}$ a constant, $\mathit{father}$ a
function of arity 1, $\mathit{older}$ a predicate of arity 2:
$\mathit{father}(\mathit{socrates})$ is a term (rule 3),
$\mathit{older}(\mathit{father}(\mathit{socrates}), X)$ is an atom (rule 4).
Exactly the same definitions hold in Prolog (Section~\ref{sec:prolog}), where variables
start with an upper-case letter, as on the board.
\end{nbwhiteboard}

\begin{nbwhy}[A small knowledge base]
\notinslides
A university information system could use: constants $\mathit{alice}$, $\mathit{bob}$,
$\mathit{sdai}$; a function $\mathit{advisor}(x)$ (the advisor of student $x$); predicates
$\mathit{student}(x)$, $\mathit{enrolled}(x, c)$, $\mathit{passed}(x, c, \mathit{grade})$.
The atom $\mathit{enrolled}(\mathit{alice}, \mathit{sdai})$ is a row in a database table;
$\mathit{enrolled}(\mathit{advisor}(\mathit{bob}), \mathit{sdai})$ says that Bob's advisor
attends SDAI. This is exactly how Prolog programs and agents' belief bases store knowledge.
\end{nbwhy}

\subsection{Quantifiers}
\label{sec:quantifiers}
\slidesrc{FOL: 14--20}

Quantification is the central part of first-order logic. With terms and predicates alone
we still cannot say:
\begin{itemize}[nosep]
  \item all men have a mother;
  \item every natural number has a prime factor.
\end{itemize}
These sentences speak about \emph{all} objects, without naming them. We need quantifiers.

\subsubsection{The three quantifiers and their syntax}
\slidesrc{FOL: 15--18}

\begin{nbdef}{Quantifiers (slide 15)}
\begin{center}
\begin{tabular}{@{}lll@{}}
\toprule
\textbf{Symbol} & \textbf{Name} & \textbf{Read} \\
\midrule
$\forall$ & universal quantifier & ``for all \ldots'' \\
$\exists$ & existential quantifier & ``there exists \ldots'' (at least one) \\
$\exists_1$ & unique quantifier & ``there exists a unique \ldots'' (exactly one) \\
\bottomrule
\end{tabular}
\end{center}
In the Z notation of the slides a quantified formula has the form
\[ \text{quantifier}\ \ \text{signature} \bullet \text{predicate} \]
where the signature is \emph{variable : type}. Example: $\forall x : \mathit{Man} \bullet
\mathit{Mortal}(x)$.
\end{nbdef}

\begin{nbexample}{reading quantified formulas (slides 17--18)}
\small
\begin{tabularx}{\linewidth}{@{}>{\raggedright\arraybackslash}p{6.2cm} X@{}}
\toprule
\textbf{Formula} & \textbf{Meaning} \\
\midrule
$\forall x : \mathit{Man} \bullet \mathit{Mortal}(x)$ & for all $x$ of type Man, $x$ is mortal (all men are mortal) \\
$\forall x : \mathit{Man} \bullet \exists_1 y : \mathit{Woman} \bullet \mathit{MotherOf}(x, y)$ & every man has exactly one mother ($y$ is the mother of $x$) \\
$\exists m : \mathit{Monitor} \bullet \mathit{MonitorState}(m, \mathit{ready})$ & there is a monitor in the ready state \\
$\forall r : \mathit{Reactor} \bullet \exists_1 t : 100..1000 \bullet \mathit{Temp}(r) = t$ & every reactor has a temperature in the range 100 to 1000 \\
$\exists n : \N \bullet n = n * n$ & some natural number is equal to its own square (true: $0$ and $1$) \\
$\exists c : \mathit{EC} \bullet \mathit{Borders}(c, \mathit{Albania})$ & some EC country borders Albania (Greece) \\
$\forall m, n : \mathit{Person} \bullet \neg \mathit{Superior}(m, n)$ & no person is superior to another \\
$\forall m : \mathit{Person} \bullet \neg \exists n : \mathit{Person} \bullet \mathit{Superior}(m, n)$ & ditto (same meaning) \\
\bottomrule
\end{tabularx}
\end{nbexample}

\paragraph{Without types.} Most AI texts (and Prolog) use untyped FOL: there is one big
domain of objects, and types become ordinary predicates. The typed formulas translate
with a fixed pattern that is worth memorising.

\begin{nbfact}{Typed to untyped quantifiers}
\notinslides
\begin{align*}
\forall x : T \bullet P(x) \quad &\text{becomes}\quad \forall x\,\big(T(x) \limp P(x)\big) \\
\exists x : T \bullet P(x) \quad &\text{becomes}\quad \exists x\,\big(T(x) \land P(x)\big)
\end{align*}
\textbf{$\forall$ goes with $\limp$, $\exists$ goes with $\land$.}
\end{nbfact}

\begin{nbwarning}[The two classic translation mistakes]
\begin{itemize}[nosep]
  \item ``All men are mortal'' is $\forall x\,(\mathit{Man}(x) \limp \mathit{Mortal}(x))$,
        \textbf{not} $\forall x\,(\mathit{Man}(x) \land \mathit{Mortal}(x))$: the latter
        says that \emph{everything} (every chair, every number) is a man and is mortal.
  \item ``Some monitor is ready'' is $\exists x\,(\mathit{Monitor}(x) \land
        \mathit{Ready}(x))$, \textbf{not} $\exists x\,(\mathit{Monitor}(x) \limp
        \mathit{Ready}(x))$: the latter is true as soon as there exists any object that is
        not a monitor (false antecedent!), even if no monitor is ready.
\end{itemize}
\end{nbwarning}

\begin{nbwarning}[The order of quantifiers matters]
\notinslides
$\forall x\,\exists y\,\mathit{MotherOf}(x, y)$: everybody has a mother (possibly a
different one for each $x$): true. \\
$\exists y\,\forall x\,\mathit{MotherOf}(x, y)$: there is one $y$ who is the mother of
everybody: false. Swapping $\forall$ and $\exists$ changes the meaning; swapping two
$\forall$ (or two $\exists$) does not.
\end{nbwarning}

\subsubsection{Free and bound variables}
\notinslides

A quantifier $\forall x$ or $\exists x$ \textbf{binds} the variable $x$ inside its
\textbf{scope} (the formula it is applied to). An occurrence of a variable that is not
bound by any quantifier is \textbf{free}. In
$\forall x\,(\mathit{gt}(x, y) \limp \mathit{gt}(x, 0))$, $x$ is bound and $y$ is free.
A formula with no free variables is a \textbf{sentence} (or \emph{closed} formula): only
sentences are simply true or false in an interpretation; a formula with free variables,
like $\mathit{gt}(x, 3)$, needs a value for its free variables (slide 13). This
distinction comes back in the proof rules of Section~\ref{sec:proof-theory-fol}.

\subsubsection{Quantifiers as big conjunctions and disjunctions}
\slidesrc{FOL: 19--20}

\begin{nbfact}{Quantifiers over a finite set (slide 19)}
Universal quantification is like conjunction, existential quantification like
disjunction:
\begin{align*}
\forall n : \{2, 4, 6\} \bullet \mathit{Even}(n) \quad &\text{is the same as}\quad
  \mathit{Even}(2) \land \mathit{Even}(4) \land \mathit{Even}(6) \\
\exists n : \{7, 8, 9\} \bullet \mathit{Prime}(n) \quad &\text{is the same as}\quad
  \mathit{Prime}(7) \lor \mathit{Prime}(8) \lor \mathit{Prime}(9)
\end{align*}
\end{nbfact}

\begin{nbfact}{Duality of $\forall$ and $\exists$ (slide 20)}
\begin{align*}
\forall x : T \bullet P(x) \quad &\liff\quad \neg\,\exists x : T \bullet \neg P(x) \\
\exists x : T \bullet P(x) \quad &\liff\quad \neg\,\forall x : T \bullet \neg P(x)
\end{align*}
Saying that everything has some property is the same as saying that there is nothing that
does not have it. Saying that something has the property is the same as saying that it is
not the case that everything lacks it.
\end{nbfact}

\begin{nbintuition}[Duality is De Morgan in disguise]
Over the finite set $\{a, b, c\}$:
\begin{align*}
\neg\,\forall x\, P(x)
  &= \neg\big(P(a) \land P(b) \land P(c)\big) && \reason{$\forall$ = big $\land$} \\
  &\equiv \neg P(a) \lor \neg P(b) \lor \neg P(c) && \reason{De Morgan} \\
  &= \exists x\, \neg P(x) && \reason{$\exists$ = big $\lor$}
\end{align*}
Negating both sides gives the first duality law. Practical reading: to refute ``all
monitors are ready'' you need one monitor that is \emph{not} ready (a counterexample).
\end{nbintuition}

\begin{nbfact}{The unique quantifier, defined with the other two}
\notinslides
$\exists_1 x\, P(x) \;\equiv\; \exists x\,\big(P(x) \land \forall y\,(P(y) \limp y = x)\big)$:
there is an $x$ with the property, and any $y$ with the property is that same $x$.
\end{nbfact}

\begin{nbexercise}{Q.1}{Translate into FOL (added)}
Use the predicates $\mathit{Student}(x)$, $\mathit{Course}(x)$, $\mathit{Attends}(x, y)$,
$\mathit{Passed}(x, y)$.
\begin{enumerate}[nosep]
  \item Every student attends some course.
  \item Some course is attended by every student.
  \item No student passed a course without attending it.
  \item Not every student passed SDAI.
\end{enumerate}
\end{nbexercise}
\begin{nbsolution}
\begin{enumerate}[nosep]
  \item $\forall x\,\big(\mathit{Student}(x) \limp \exists y\,(\mathit{Course}(y) \land
        \mathit{Attends}(x, y))\big)$
  \item $\exists y\,\big(\mathit{Course}(y) \land \forall x\,(\mathit{Student}(x) \limp
        \mathit{Attends}(x, y))\big)$ \ (stronger than 1: the order of quantifiers changed)
  \item $\neg \exists x\, \exists y\, \big(\mathit{Student}(x) \land \mathit{Passed}(x, y)
        \land \neg \mathit{Attends}(x, y)\big)$, equivalently
        $\forall x\, \forall y\, \big(\mathit{Student}(x) \land \mathit{Passed}(x, y)
        \limp \mathit{Attends}(x, y)\big)$ (by duality and De Morgan).
  \item $\neg \forall x\,(\mathit{Student}(x) \limp \mathit{Passed}(x, \mathit{sdai}))$,
        equivalently $\exists x\,(\mathit{Student}(x) \land \neg \mathit{Passed}(x,
        \mathit{sdai}))$: some student did not pass SDAI.
\end{enumerate}
\textbf{Lesson}: always check the pattern $\forall \ldots \limp$, $\exists \ldots \land$.
\end{nbsolution}

\subsection{Decidability}
\label{sec:decidability}
\slidesrc{FOL: 21}

In propositional logic some formulas are tautologies (true under all interpretations), and
there is a \emph{procedure}, the truth-table method, that tells in a finite time whether any
formula is a tautology.

\begin{nbdef}{Valid formula (slide 21)}
A formula of FOL that is true under all interpretations is said to be \textbf{valid}
(it is the FOL analogue of a tautology).
\end{nbdef}

Truth tables do not work for FOL: an interpretation chooses a domain of objects (possibly
infinite, like $\N$) and a meaning for every symbol, so there are infinitely many
interpretations and a ``truth table'' would be infinite. Is there \emph{any other}
procedure guaranteed to tell us, in a finite amount of time, whether a FOL formula is or
is not valid? \textbf{No.}

\begin{nbfact}{Undecidability of FOL (Church and Turing, 1936)}
There is no algorithm that, given any FOL formula, always terminates and correctly answers
whether the formula is valid. FOL is said to be \textbf{undecidable}.
\end{nbfact}

\begin{nbintuition}[Semi-decidable: ``yes'' can be found, ``no'' may never come]
\notinslides
FOL is \emph{semi-decidable}: there are procedures (e.g.\ enumerating all proofs, or
resolution) that, if the formula \emph{is} valid, will find a proof in finite time. If the
formula is \emph{not} valid, the procedure may run forever, and at no point can you tell
``not valid'' apart from ``valid, but the proof has not been found yet''. You will meet
exactly this behaviour in Prolog: a query either succeeds, fails, or loops forever
(Section~\ref{sec:prolog}).
\end{nbintuition}

\begin{center}
\begin{tikzpicture}[font=\small, >=Stealth,
  b/.style={draw, rounded corners, thick, align=center, text width=4.6cm, minimum height=1.3cm}]
  \node[b, draw=nbgreen, fill=nbgreen!8] (p) at (0,0) {\textbf{Propositional logic}\\decidable: truth tables always terminate (but $2^n$ rows)};
  \node[b, draw=nbred, fill=nbred!8] (f) at (6.4,0) {\textbf{First-order logic}\\undecidable; semi-decidable (valid formulas have finite proofs)};
  \draw[->, thick] (p) -- node[above, font=\scriptsize]{more expressive} node[below, font=\scriptsize]{less computable} (f);
\end{tikzpicture}
\captionof{figure}{The price of expressive power: FOL can say much more, but no algorithm can
decide validity for all its formulas.}
\end{center}

\begin{nbwhy}
This is why practical symbolic AI uses \emph{fragments} of FOL with better computational
behaviour: Prolog uses Horn clauses with a specific search strategy
(Section~\ref{sec:prolog}), description logics (behind ontologies and the Semantic Web) are
decidable fragments, and agent languages such as AgentSpeak keep beliefs as ground atoms.
\end{nbwhy}

\subsection{Proof theory for first-order logic}
\label{sec:proof-theory-fol}
\slidesrc{Proof theory (FOL): 1--4}

The source for this part is a set of university lecture notes (Math 260A, UCSD, 2012),
more technical than the rest of the course. The goal here is to understand the
\emph{ideas}: what a proof system for FOL looks like, what soundness means, and how
quantifiers are handled.

In propositional logic the simplest proof system is the truth table. In FOL there is no
analogue: the truth tables would be infinite (Section~\ref{sec:decidability}). We go
directly to \textbf{Frege-style} systems, called \textbf{Hilbert-style} systems in FOL:
\emph{many axioms, very few rules}.

\subsubsection{Frege (Hilbert-style) proofs}
\slidesrc{Proof theory (FOL): 1--2}

\begin{nbdef}{Axioms of the Frege system for FOL (page 1)}
\begin{itemize}
  \item \textbf{Propositional axioms}: all axioms of propositional Frege systems, such as
        $\varphi \limp (\psi \limp \varphi)$, where now $\varphi, \psi$ can be any FOL
        formulas.
  \item \textbf{Quantifier axioms}: for any formula $\varphi$ and any term $t$,
        \[ \varphi(t) \limp \exists x\,\varphi(x) \qquad\qquad \forall x\,\varphi(x) \limp \varphi(t) \]
        (provided $t$ is \emph{substitutable} for $x$; the notes mostly ignore this
        technical condition).
  \item \textbf{Equality axioms} (if $=$ is in the language):
        \begin{itemize}[nosep]
          \item $\forall x\,(x = x)$ \hfill(reflexivity)
          \item $\forall x \forall y\,(x = y \limp y = x)$ \hfill(symmetry)
          \item $\forall x \forall y \forall z\,(x = y \land y = z \limp x = z)$ \hfill(transitivity)
          \item $\forall \vec x\, \forall \vec y\,\big(x_1 = y_1 \land \cdots \land x_k = y_k
                \limp f(x_1, \ldots, x_k) = f(y_1, \ldots, y_k)\big)$ for every function $f$
                of arity $k$
          \item $\forall \vec x\, \forall \vec y\,\big(x_1 = y_1 \land \cdots \land x_k = y_k
                \limp (P(x_1, \ldots, x_k) \limp P(y_1, \ldots, y_k))\big)$ for every
                relation $P$ of arity $k$ \hfill(substitution of equals)
        \end{itemize}
\end{itemize}
\end{nbdef}

\begin{nbwarning}[Corrections to page 1 of the notes]
\begin{itemize}[nosep]
  \item The notes say ``$t$ should be substitutable for $t$'': it should read ``for $x$''.
  \item The last two equality axioms are printed without the hypothesis
        $x_1 = y_1 \land \cdots \land x_k = y_k$, which would make them say that $f$ gives
        the same value on \emph{any} arguments (false). The correct versions are written
        above: equal inputs give equal outputs, and equal objects have the same properties.
\end{itemize}
\end{nbwarning}

\begin{nbintuition}[Reading the quantifier axioms]
$\forall x\,\varphi(x) \limp \varphi(t)$: ``if something holds for everything, it holds for
any particular thing $t$'' (\emph{instantiation}). From ``all humans are mortal'' you may
conclude ``if Socrates is human, Socrates is mortal''.
$\varphi(t) \limp \exists x\,\varphi(x)$: ``if it holds for some particular $t$, then it
holds for something'' (\emph{generalisation to $\exists$}). From
$\mathit{Ready}(X12)$ conclude ``some monitor is ready''.
\end{nbintuition}

\begin{nbfact}{Rules of inference (pages 1--2)}
\[
\text{Modus Ponens:}\ \ \frac{\varphi \qquad \varphi \limp \psi}{\psi}
\qquad\qquad
\frac{\psi \limp \varphi}{\psi \limp \forall x\,\varphi}
\qquad\qquad
\frac{\varphi \limp \psi}{\exists x\,\varphi \limp \psi}
\]
The two quantifier rules require that \textbf{$x$ is not free in $\psi$}.
\end{nbfact}

\begin{nbintuition}[Why ``$x$ not free in $\psi$''?]
Rule 2: if $\psi \limp \varphi(x)$ has been proved and $\psi$ says nothing about $x$, then
the proof worked for an arbitrary $x$, so $\psi \limp \forall x\,\varphi$. Without the
condition we could go from $x = 0 \limp x = 0$ to $x = 0 \limp \forall x\,(x = 0)$, i.e.\
``if this number is 0 then every number is 0'': wrong. Rule 3: if $\varphi(x)$ implies
$\psi$ whatever $x$ is, then knowing that \emph{some} $x$ satisfies $\varphi$ is enough
to get $\psi$, provided $\psi$ itself does not depend on which $x$ it was.
\end{nbintuition}

\begin{nbwhiteboard}[label=wb:modus-ponens]{Modus Ponens: Socrates is mortal}{written during the FOL proof-theory lecture}
\begin{center}
\begin{tabular}{rl}
\textcircled{\scriptsize 1} & $\forall x\,.\ \mathit{human}(x) \Rightarrow \mathit{mortal}(x)$ \\
\textcircled{\scriptsize 2} & $\mathit{human}(\mathit{socrates})$ \\
\hline
 & $\mathit{mortal}(\mathit{socrates})$ \qquad \textsc{Modus Ponens}
\end{tabular}
\end{center}
\wbphoto[0.6\linewidth]{sections/first-order-logic/proof-theory-fol/assets/whiteboard-6-2.jpg}
The board shows Aristotle's syllogism (Section~\ref{sec:ai-history}) in FOL. Strictly
speaking, Modus Ponens needs an implication whose antecedent is \emph{exactly}
$\mathit{human}(\mathit{socrates})$, so a quantifier step comes first. In full, with the
Frege system above:
Writing $h$ for $\mathit{human}$, $m$ for $\mathit{mortal}$ and $s$ for
$\mathit{socrates}$:
\begin{center}
\small
\begin{tabular}{rll}
1 & $\forall x\,(h(x) \limp m(x))$ & premise \\
2 & $h(s)$ & premise \\
3 & $\forall x\,(h(x) \limp m(x)) \limp (h(s) \limp m(s))$ & quantifier axiom with $t = s$ \\
4 & $h(s) \limp m(s)$ & Modus Ponens 1, 3 \\
5 & $m(s)$ & Modus Ponens 2, 4
\end{tabular}
\end{center}
In Prolog the same knowledge is \code{mortal(X) :- human(X).} and \code{human(socrates).};
the query \code{?- mortal(socrates).} answers \code{true}: Prolog's resolution does the
instantiation ($X = \mathit{socrates}$, called \emph{unification}) and the Modus Ponens
step in one go (Section~\ref{sec:prolog}).
\end{nbwhiteboard}

\subsubsection{Provability and soundness}
\slidesrc{Proof theory (FOL): 2}

\begin{nbdef}{Provability $\Gamma \proves \varphi$ (page 2)}
Let $\Gamma$ be a set of formulas (the hypotheses). We write $\Gamma \proves \varphi$
(``$\Gamma$ proves $\varphi$'') if there is a finite sequence of formulas, each coming from
$\Gamma$, from the axioms, or from earlier formulas by an inference rule, ending at
$\varphi$.
\end{nbdef}

Recall that $\Gamma \entails \varphi$ (``$\Gamma$ entails $\varphi$'') means: in every
structure (interpretation) $M$ in which all formulas of $\Gamma$ are true, $\varphi$ is
true too. $\proves$ is about \emph{symbols and rules}; $\entails$ is about \emph{meaning and
truth}.

\begin{nbfact}{Soundness theorem (page 2)}
If $\Gamma \proves \varphi$, then $\Gamma \entails \varphi$. \quad (Everything provable is true.)
\end{nbfact}

\paragraph{Proof idea.} Take a proof $P$ of $\varphi$ from $\Gamma$ and a structure $M$ in
which $\Gamma$ is true, i.e.\ for all object assignments $\sigma$ (values for the free
variables) and all $\gamma \in \Gamma$, $M \entails \gamma[\sigma]$. Show by
\textbf{induction on the number of lines} of $P$ that every line $\psi$ is true in $M$
under all $\sigma$. Base cases: hypotheses (true by assumption) and axioms (true in every
structure). Modus Ponens preserves truth (truth table of $\limp$). For the
$\exists$-rule $\frac{\varphi \limp \psi}{\exists x\varphi \limp \psi}$ argue by
contrapositive:
\begin{enumerate}[nosep]
  \item suppose $M \nentails (\exists x\varphi \limp \psi)[\sigma]$; then
        $M \entails \exists x\varphi[\sigma]$ and $M \nentails \psi[\sigma]$;
  \item the first fact gives an object $m$ of $M$ with $M \entails \varphi[\sigma(m/x)]$
        ($\sigma$ changed so that $x$ denotes $m$);
  \item let $\sigma' = \sigma(m/x)$. Since $x$ is not free in $\psi$, changing the value of
        $x$ does not affect $\psi$, so $M \nentails \psi[\sigma']$;
  \item under $\sigma'$, $\varphi$ is true and $\psi$ false, so
        $M \nentails (\varphi \limp \psi)[\sigma']$: the premise of the rule is not true
        under all assignments. Contrapositive: if the premise is, so is the conclusion.
\end{enumerate}
(The notes write $M \entails \psi[\sigma(m/x)]$ in step 2: it should be $\varphi$.) Note
that we needed to \emph{change} the object assignment to make this work.

\begin{nbfact}{Completeness (G\"odel, 1929)}
\notinslides
The converse also holds: if $\Gamma \entails \varphi$ then $\Gamma \proves \varphi$. So
for FOL, ``provable'' and ``true in all models'' coincide. This does not contradict
undecidability: a proof always \emph{exists} for valid formulas, but there is no
algorithm that also tells you, in finite time, when \emph{no} proof exists.
\end{nbfact}

\begin{center}
\begin{tikzpicture}[font=\small, >=Stealth]
  \node[draw, rounded corners, thick, nbpurple, fill=nbpurple!6, align=center, minimum width=3.5cm, minimum height=1.1cm] (s) at (0,0) {$\Gamma \proves \varphi$\\\footnotesize syntax: there is a proof};
  \node[draw, rounded corners, thick, nbgreen, fill=nbgreen!6, align=center, minimum width=3.5cm, minimum height=1.1cm] (m) at (6.5,0) {$\Gamma \entails \varphi$\\\footnotesize semantics: true in every model};
  \draw[->, thick] (s.north east) to[bend left=15] node[above]{soundness} (m.north west);
  \draw[->, thick] (m.south west) to[bend left=15] node[below]{completeness} (s.south east);
\end{tikzpicture}
\captionof{figure}{Soundness: the rules never prove something false. Completeness: every
truth that follows has a proof.}
\end{center}

\subsubsection{Sequent calculus LK}
\slidesrc{Proof theory (FOL): 2--4}

A second style of proof system, invented by Gentzen: the \textbf{sequent calculus}. A
\textbf{sequent} has the form
\[ \Gamma \to \Delta \qquad (\Gamma, \Delta \text{ finite lists of formulas}) \]
and means: \emph{if all formulas in $\Gamma$ are true, then at least one formula in
$\Delta$ is true}. LK (from the German \emph{Logischer Kalk\"ul}) has the opposite
philosophy of Frege systems: \emph{one axiom, many rules}.

\paragraph{Variable conventions.} $a, b, c, \ldots$ are used only for \emph{free}
variables and $x, y, z, \ldots$ only for \emph{bound} variables. Consequence: a
sub-formula of a formula is not always a formula. In $\forall x\,(x \leq a)$, the part
$x \leq a$ is a \emph{semi-formula} ($x$ is not bound there); similarly $x + a$ is a
\emph{semi-term}. With these conventions any term $t$ is always substitutable for a free
variable $a$ in a formula $\varphi$.

\begin{nbfact}{LK (page 3)}
\textbf{Initial sequents}: $A \to A$, for any formula $A$. \\
\textbf{Rules}: all rules of the propositional sequent calculus, plus four quantifier rules:
\begin{center}
\begin{tabular}{cc}
\AxiomC{$\Gamma \to \Delta, A(b)$}
\LeftLabel{\small $\forall$-right}
\UnaryInfC{$\Gamma \to \Delta, \forall x A(x)$}
\DisplayProof
&
\AxiomC{$\Gamma \to \Delta, A(t)$}
\LeftLabel{\small $\exists$-right}
\UnaryInfC{$\Gamma \to \Delta, \exists x A(x)$}
\DisplayProof
\\[18pt]
\AxiomC{$A(b), \Gamma \to \Delta$}
\LeftLabel{\small $\exists$-left}
\UnaryInfC{$\exists x A(x), \Gamma \to \Delta$}
\DisplayProof
&
\AxiomC{$A(t), \Gamma \to \Delta$}
\LeftLabel{\small $\forall$-left}
\UnaryInfC{$\forall x A(x), \Gamma \to \Delta$}
\DisplayProof
\end{tabular}
\end{center}
In $\forall$-right and $\exists$-left the free variable $b$ is an \textbf{eigenvariable}:
it must not appear anywhere in the lower sequent (the substitution $x/b$ ``captures'' all
occurrences of $b$). In $\exists$-right and $\forall$-left, $t$ is \emph{any} term and may
occur in the lower sequent.
\end{nbfact}

\begin{nbintuition}[Eigenvariable = ``an arbitrary object we know nothing about'']
\begin{itemize}[nosep]
  \item \emph{$\forall$-right}: to prove ``for all $x$, $A(x)$'', prove $A(b)$ for a
        \emph{fresh} $b$ about which the hypotheses say nothing. Like in maths: ``let $b$
        be an arbitrary natural number\ldots''. If $b$ appeared in $\Gamma$ (e.g.\
        $\Gamma$ says $b = 0$), it would not be arbitrary.
  \item \emph{$\exists$-left}: if we know ``some $x$ has $A$'', we may give it a fresh name
        $b$, but we cannot assume it is any object we already know.
  \item \emph{$\exists$-right} and \emph{$\forall$-left}: to show ``something has $A$'',
        exhibit any witness $t$; to use ``everything has $A$'', instantiate with any $t$.
\end{itemize}
\end{nbintuition}

\subsubsection{LKe: adding equality}
\slidesrc{Proof theory (FOL): 4}

LK has no rules for $=$. \textbf{LKe} adds five initial sequents (for any terms
$s, t, u, s_i, t_i$):
\begin{itemize}[nosep]
  \item $\to s = s$
  \item $s = t \to t = s$
  \item $s = t,\ t = u \to s = u$
  \item $s_1 = t_1, \ldots, s_k = t_k \to f(s_1, \ldots, s_k) = f(t_1, \ldots, t_k)$ for any
        \emph{function symbol} $f$ of arity $k$
  \item $s_1 = t_1, \ldots, s_k = t_k,\ P(s_1, \ldots, s_k) \to P(t_1, \ldots, t_k)$ for
        any relation $P$ of arity $k$
\end{itemize}
(The notes write ``for any formula $f$'': $f$ is a function symbol. The last sequent is
printed with a second arrow; the version above is the intended one.)

\begin{nbexample}{proving symmetry of equality in LKe (page 4)}
Goal: $\to \forall x \forall y\,(x = y \limp y = x)$. Read the proof bottom-up as a plan,
top-down as a derivation:
\begin{center}
\AxiomC{$a = b \to b = a$}
\RightLabel{\small $\limp$-right}
\UnaryInfC{$\to a = b \limp b = a$}
\RightLabel{\small $\forall$-right ($b$ fresh)}
\UnaryInfC{$\to \forall y\,(a = y \limp y = a)$}
\RightLabel{\small $\forall$-right ($a$ fresh)}
\UnaryInfC{$\to \forall x \forall y\,(x = y \limp y = x)$}
\DisplayProof
\end{center}
\begin{enumerate}[nosep]
  \item Start from the LKe initial sequent $a = b \to b = a$ (symmetry for the terms $a, b$).
  \item $\limp$-right (a propositional rule: from $A, \Gamma \to \Delta, B$ infer
        $\Gamma \to \Delta, A \limp B$) moves $a = b$ to the right as an antecedent.
  \item $\forall$-right with eigenvariable $b$: $b$ does not occur in the lower sequent
        $\to \forall y\,(a = y \limp y = a)$, so the condition holds.
  \item $\forall$-right again with eigenvariable $a$.
\end{enumerate}
\end{nbexample}

\subsubsection{Natural deduction rules for quantifiers}
\notinslides

To connect with the natural deduction of Section~\ref{sec:natural-deduction-prop}, here are
the four quantifier rules in the same Intro/Elim style (same ideas as LK):
\begin{center}
\small
\renewcommand{\arraystretch}{1.3}
\begin{tabularx}{\linewidth}{@{}l X@{}}
\toprule
$\forall$\,Elim & from $\forall x\,\varphi(x)$ write $\varphi(t)$ for any term $t$ (instantiation) \\
$\forall$\,Intro & if you derived $\varphi(b)$ for a fresh $b$ appearing in no open
                   assumption, write $\forall x\,\varphi(x)$ \\
$\exists$\,Intro & from $\varphi(t)$ write $\exists x\,\varphi(x)$ \\
$\exists$\,Elim & from $\exists x\,\varphi(x)$, assume $\varphi(b)$ with $b$ fresh in a
                  subproof; if you derive $\psi$ (not containing $b$), write $\psi$ outside \\
\bottomrule
\end{tabularx}
\end{center}

\begin{nbexample}{the monitor argument, proved}
\begin{fitchproof}
\hypo{1}{\forall x\,(\mathit{Monitor}(x) \limp \mathit{Ready}(x))} \RPrem
\hypo{2}{\mathit{Monitor}(X12)} \RPrem
\have{3}{\mathit{Monitor}(X12) \limp \mathit{Ready}(X12)} \by{$\forall$\,Elim}{1}
\have{4}{\mathit{Ready}(X12)} \RImpE{3,2}
\end{fitchproof}
The argument of Section~\ref{sec:why-fol}, which propositional logic could not capture, becomes a two-step proof.
\end{nbexample}


\begin{nbsummary}[First-order logic]
\begin{itemize}[nosep]
  \item \textbf{Terms} denote objects: constants ($\mathit{socrates}$, $8$), variables
        ($x$), functional terms $f(t_1,\ldots,t_n)$ ($\mathit{plus}(2,3)$,
        $\mathit{mother}(x)$). A term is never true or false.
  \item \textbf{Atoms} (predicates): $P(t_1,\ldots,t_n)$, true or false under an
        interpretation. Arity 1 = \emph{property}; arity 0 = \emph{primitive proposition};
        an atom without variables is a \emph{proposition}.
  \item Formulas: atoms combined with $\neg, \land, \lor, \limp, \liff$ and quantifiers
        $\forall x$ (for all), $\exists x$ (there exists), $\exists_1 x$ (exactly one).
  \item Translation pattern: ``all $A$ are $B$'' = $\forall x\,(A(x) \limp B(x))$;
        ``some $A$ is $B$'' = $\exists x\,(A(x) \land B(x))$.
  \item $\forall$ over a finite set = big $\land$; $\exists$ = big $\lor$. Duality:
        $\forall x\,P(x) \equiv \neg\exists x\,\neg P(x)$,
        $\exists x\,P(x) \equiv \neg\forall x\,\neg P(x)$.
  \item \textbf{Valid} = true under all interpretations. Validity in FOL is
        \textbf{undecidable} (no algorithm always answers yes/no).
  \item $\Gamma \proves \varphi$: $\varphi$ is provable from $\Gamma$ (syntax, rules).
        $\Gamma \entails \varphi$: $\varphi$ is true in every model of $\Gamma$ (semantics).
        \textbf{Soundness}: $\Gamma \proves \varphi \Rightarrow \Gamma \entails \varphi$.
  \item Frege/Hilbert system: many axioms + Modus Ponens + two quantifier rules.
        Sequent calculus LK: one axiom $A \to A$ + many rules; $\forall$-right and
        $\exists$-left need a fresh \emph{eigenvariable}.
  \item Socrates (Whiteboard~\ref{wb:modus-ponens}): instantiate
        $\forall x(\mathit{human}(x) \limp \mathit{mortal}(x))$ with $x = \mathit{socrates}$,
        then Modus Ponens.
\end{itemize}
\end{nbsummary}
